% ============================================================
% Chapter 6 --- Completeness and the Banach fixed point theorem
% Originally: content/week-02.tex, \section{Completeness} (Corsin Week 2, Friday)
%           + content/week-03.tex, \subsection{Contraction mappings and the Banach fixed
%             point theorem} (Corsin Week 3, Friday)
%
% Placed after \cref{ch:continuity} because a contraction is by definition a Lipschitz map
% with constant < 1, so \cref{def:lipschitz_continuous} must already exist. Corsin's own
% order was the same; only the week boundary between the two halves is gone.
% ============================================================
\chapter{Completeness and the Banach Fixed Point Theorem}
\label{ch:completeness}

% Transition: Claude Opus 5 (High)
Having examined how maps behave between metric spaces, we return to a property of a single space.
Does every sequence that \qt{looks like} it ought to converge (\cref{def:convergent_sequence})
actually do so? A space where the answer is yes is called complete.

Completeness is the first genuinely \emph{metric} rather than topological property met in this
part, and the distinction is sharp: two metrics can generate exactly the same open sets while
only one of them is complete. What completeness is usually good for is producing objects one
cannot write down. The Banach fixed point theorem at the end of the chapter is the standard
instrument for that, and it reappears as the engine of both the inverse function theorem and the
Picard--Lindelöf existence theorem later on.

% Originally: content/week-02.tex, \section{Completeness} (Corsin Week 2, Friday)

\section{Completeness}
\label{sec:completeness}

\begin{definition}[Cauchy sequence and complete metric space]
\label{def:complete_metric_space}
Let $(X,d)$ be a metric space and $(x_n)_{n=0}^{\infty}$ a sequence in $X$. We say that
$(x_n)_{n=0}^{\infty}$ is \newterm{Cauchy} if
\[\forall \varepsilon > 0 \ \exists N \in \mathbb{N} \text{ such that } d(x_n, x_m) < \varepsilon \quad \forall n, m \geq N.\]
$(X,d)$ is a \newterm{complete metric space} (\germanterm{vollst\"andiger metrischer Raum}) if
every Cauchy sequence in $X$ converges with respect to $d$.
\end{definition}

% Supplement: Sascha Brack/Class Notes/Week_02_Notes_Monday_Updated.pdf, pp. 7--9
\begin{proposition}[Convergence in $\mathbb{R}^n$ is coordinatewise; $\mathbb{R}^n$ is complete]
\label{prop:Rn_coordinatewise_complete}
Equip $\mathbb{R}^n$ with the standard metric $d_2$, and write
$x_m = (x_{m,1},\dots,x_{m,n})$.
\begin{enumerate}[label=\textbf{(\alph*)}]
  \item $x_m \to x$ in $\mathbb{R}^n$ if and only if $x_{m,i} \to x_i$ in $\mathbb{R}$ for every
    $i = 1,\dots,n$.
  \item $(\mathbb{R}^n, d_2)$ is complete.
\end{enumerate}
\end{proposition}

\begin{proof}
Everything follows from the two-sided estimate, valid for every $j$:
\[|x_{m,j}-x_j| \ \leq\ \lVert x_m-x\rVert \ =\ \sqrt{\sum_{i=1}^n|x_{m,i}-x_i|^2}
  \ \leq\ \sqrt{n}\,\max_i|x_{m,i}-x_i|.\]

\textbf{(a)} If $\lVert x_m-x\rVert \to 0$, the left inequality forces each coordinate to
converge. Conversely, if every coordinate converges then $\max_i|x_{m,i}-x_i| \to 0$ (a maximum
of \emph{finitely} many null sequences), so the right inequality forces
$\lVert x_m - x\rVert \to 0$.

\textbf{(b)} Let $(x_m)$ be Cauchy in $\mathbb{R}^n$. The same left inequality gives
$|x_{m,j}-x_{k,j}| \leq \lVert x_m-x_k\rVert$, so each coordinate sequence $(x_{m,j})_m$ is
Cauchy in $\mathbb{R}$. Since $\mathbb{R}$ is complete --- this is the completeness axiom from
\emph{Analysis I}, and the only input the proof has --- each converges, say
$x_{m,j} \to x_j$. Setting $x := (x_1,\dots,x_n)$, part \textbf{(a)} gives $x_m \to x$.
\end{proof}

\begin{ainote}
Worth noting how little is proved here: completeness of $\mathbb{R}^n$ is completeness of
$\mathbb{R}$, applied $n$ times. That is also where the argument would fail in infinite
dimensions --- $\max_i$ over infinitely many coordinates need not tend to $0$ even when each
does, which is exactly the phenomenon behind the $\varepsilon$-net argument in
\Cref{lem:sequentially_compact_totally_bounded}.
\end{ainote}

% Source: Corsin Nick/Class Notes/Week 2.pdf, pp. 1--13
\begin{example}[Completeness examples]
\begin{itemize}
  \item The Euclidean spaces $(\mathbb{R}^n, \langle\cdot,\cdot\rangle_{\text{eucl}})$ are
    \textbf{complete}.
  \item Any \textbf{non-closed} subset $U \subseteq \mathbb{R}^n$ (standard metric) and the
    rationals $\mathbb{Q}$ are \textbf{not complete}.
  \item $\big(C^0([0,1]), \sup_{x \in [0,1]}|\cdot|\big)$ is \textbf{complete} (Zorich, p.\ 22).
\end{itemize}
\end{example}

% Generator: Claude Opus 5 (Medium)
\begin{remark}[Completeness is not the same as closedness]
\label{rem:complete_vs_closed}
The second bullet is easy to over-read as \qt{complete means closed}. It does not, and the
difference is worth pinning down now (because it recurs).

\textbf{Closedness is relative, completeness is intrinsic.} \qt{Closed} only means anything once
you say \emph{closed in what} (i.e., it is a statement about how a set sits inside an ambient space).
Completeness asks nothing about an ambient space --- only whether the Cauchy sequences of the
space converge \emph{within it}. So:
\begin{itemize}
  \item $\mathbb{Q}$ is \textbf{closed in $\mathbb{Q}$} --- every space is closed in itself
    (\cref{ex:X_empty_clopen}) --- and yet \textbf{not complete}: the sequence
    $1,\,1.4,\,1.41,\,1.414,\dots$ is Cauchy in $\mathbb{Q}$ and converges to nothing there.
  \item $(0,1)$ is \textbf{not closed in $\mathbb{R}$} and also not complete, which is the
    second bullet's case --- but the two failures have different causes, and the first does not
    imply the second in general.
\end{itemize}

What \emph{is} true is a one-directional statement, and only inside a complete ambient space:
if $X$ is complete and $A \subseteq X$, then
\[A \text{ closed in } X \iff A \text{ complete}.\]
This is why the second bullet is correct \emph{as stated for $\mathbb{R}^n$} --- it silently uses
completeness of the ambient $\mathbb{R}^n$. Drop that hypothesis and the equivalence dies, as
$\mathbb{Q}$ inside $\mathbb{Q}$ shows.

\begin{ainote}
The same relative-versus-intrinsic distinction is the one behind
\Cref{rem:common_misconception} ($[0,1]^2$ is open in itself, not in $\mathbb{R}^2$), and it
returns as \cref{rem:boundedness_not_topological}. Compactness, by contrast, is
intrinsic like completeness: a subset $K$ is compact exactly when the metric space
$(K, d|_{K\times K})$ is, with no reference to the ambient space at all --- which is precisely
the reduction used in \cref{def:sequential_topological_compactness}.
\end{ainote}
\end{remark}

% Quelle: Diego Torres Tejeda/Notes - 06.03 - Diego Torres Tejeda.pdf, p. 2
\begin{exercise}[Completeness is not a topological property]
\label{ex:completeness_not_topological}
Corsin does not address whether completeness survives a change of metric; Diego Torres Tejeda
sets this exercise, which settles it. Let $X := (0,1]$, let $d_{\text{eucl}}(x,y) := |x-y|$ be
the standard metric on $X$, and define
\[d(x,y) := \left|\frac1x - \frac1y\right|, \qquad x,y \in X.\]
\begin{enumerate}[label=\textbf{(\alph*)}]
  \item Show that $d$ is a metric on $X$.
  \item Show that $(X,d_{\text{eucl}})$ is \textbf{not} complete.
  \item Show that $(X,d)$ \textbf{is} complete.
  \item Show that $d_{\text{eucl}}$ and $d$ induce the \textbf{same topology} on $X$.
\end{enumerate}
\textit{Diego's hint for} \textbf{(c)}: let $(x_n)$ be $d$-Cauchy. Then $(1/x_n)$ is Cauchy for
the standard metric, so it has a limit $\ell \in \mathbb{R}$; show $\ell \geq 1$ and that
$x_n \to 1/\ell$.
\end{exercise}

\begin{ainote}
Diego states the exercise on $X = (0,\infty)$, but there part \textbf{(c)} is false: on
$(0,\infty)$ the map $x\mapsto1/x$ is a bijection \emph{onto} $(0,\infty)$, so $(X,d)$ is
isometric to $\big((0,\infty),|\cdot|\big)$, which is not complete --- concretely $x_n := n$
satisfies $d(x_n,x_m) = |\tfrac1n-\tfrac1m| \to 0$ yet has no limit in $X$. His own hint asks to
show $\ell > 0$, and that is precisely the step that fails. Restricting to a bounded interval
repairs it: on $(0,1]$ the image is $[1,\infty)$, which \emph{is} closed in $\mathbb{R}$ and
hence complete. Typeset with $X := (0,1]$ above.
\end{ainote}

% Sonnet 5 (Medium)
Compactness will turn out to be a strictly stronger finiteness-type property than completeness:
every compact metric space is complete (\cref{ex:compact_finite_complete}), but the converse
fails --- $\mathbb{R}^n$ itself is complete, yet (being unbounded) not compact, as we will see
next (\cref{thm:sequential_compactness}, \cref{sec:heine_borel}).

% Source: Corsin Nick/Class Notes/Week 2.pdf, pp. 1--13

% Supplement: Analysis_II_Script_v1.pdf, pp. 11--13
\section{The completion of a metric space}
\label{sec:completion_of_a_metric_space}

% This section has no counterpart in Corsin's notes. It is transcribed directly from the official
% script (S9.1.3, printed pp. 11--13), rather than enriching an existing tutor passage the way the
% rest of this document's supplements do. Included because script Exercises 9.33--9.35 walking
% through the construction are short, self-contained, and exactly the kind of material this
% chapter collects. Not included: the script's own use of the construction to build R from Q
% (S9.1.4), which it marks "extra; cf. Grundstrukturen" and which belongs to the sibling project.
% Transition: Claude Opus 5 (High)
\Cref{rem:complete_vs_closed} left one question open. We know $\mathbb{Q}$ is not complete and
$\mathbb{R}$ is, and that $\mathbb{R}$ is built to repair exactly that defect. Is the repair a
one-off trick special to the reals, or can any incomplete space be fixed the same way?

It can, and the construction below is the general answer. Every metric space embeds isometrically,
and densely, into a complete one, obtained by taking the space of its own Cauchy sequences and
declaring two of them equal when they ought to have the same limit. The reals from the rationals
is then just the first instance.

\begin{definition}[Completion of a metric space]
\label{def:completion_of_metric_space}
Let $(X,d)$ be a metric space. Write $\mathcal{C}_X$ for the set of all Cauchy sequences in $X$,
and define an equivalence relation on $\mathcal{C}_X$ by
\[(x_n)_{n=0}^\infty \sim (y_n)_{n=0}^\infty \quad :\Longleftrightarrow\quad \lim_{n\to\infty}d(x_n,y_n) = 0.\]
The quotient set $\overline X := \mathcal{C}_X/{\sim}$, equipped with the map
\[\overline d\big([(x_n)_{n=0}^\infty],\,[(y_n)_{n=0}^\infty]\big) := \lim_{n\to\infty}d(x_n,y_n),\]
is called the \newterm{completion} (\germanterm{Vervollst\"andigung}) of $(X,d)$. The injection
$\iota : X \to \overline X$ sending $x$ to the class of the constant sequence $(x,x,x,\dots)$ is
called the \newterm{canonical embedding}. For all $x,y \in X$,
\[d(x,y) = \overline d\big(\iota(x),\iota(y)\big),\]
which in particular shows $\iota$ is injective.
\end{definition}

% Supplement: Analysis_II_Script_v1.pdf, p. 11
\begin{exercise}[Well-definedness of the completion]
\label{ex:completion_well_defined}
Show that the objects introduced in \cref{def:completion_of_metric_space} are well defined. In
particular, verify that $\overline d$ is indeed a metric on $\overline X$.
\end{exercise}

% Supplement: Analysis_II_Script_v1.pdf, p. 11
\begin{exercise}[The completion is complete]
\label{ex:completion_is_complete}
As the name suggests, the completion $(\overline X,\overline d)$ of a metric space is complete,
meaning every Cauchy sequence in $\overline X$ converges. A sequence in $\overline X$ is
essentially a sequence of sequences, i.e.\
$\big[(x_{m,n})_{n=0}^\infty\big]_{m=0}^\infty$. Show that $\big[(x_{n,n})_{n=0}^\infty\big]$ is a
limit of this sequence.
\end{exercise}

% Supplement: Analysis_II_Script_v1.pdf, p. 11
\begin{exercise}[Universal property of the completion]
\label{ex:completion_universal_property}
Let $(X,d)$ be a metric space with completion $(\overline X,\overline d)$. Let $(Y,d_Y)$ be a
complete metric space, and let $f : X \to Y$ be a function such that
\[d(x,y) = d_Y\big(f(x),f(y)\big) \qquad \text{for all } x,y \in X.\]
Show that there exists a unique function $\overline f : \overline X \to Y$ such that
$f = \overline f \circ \iota_X$ and
\[\overline d(\overline x,\overline y) = d_Y\big(\overline f(\overline x),\overline f(\overline y)\big)
  \qquad \text{for all } \overline x,\overline y \in \overline X.\]
This can be interpreted as: \qt{$\overline X$ is the \emph{smallest} complete metric space
containing $X$.}
\end{exercise}

\begin{remark}
\Cref{ex:completion_universal_property} is what makes \qt{the} completion well-posed rather than
just \qt{a} completion: any other complete space that $X$ embeds into isometrically receives a
unique isometric copy of $\overline X$ as well, via $\overline f$. This is the same
uniqueness-up-to-isomorphism pattern behind \qt{the} real numbers as \qt{the} complete ordered
field, and it is exactly how the script uses this construction in $\S9.1.4$, taking
$X = \mathbb{Q}$.
\end{remark}

% Originally: content/week-03.tex, \subsection{Contraction mappings and the Banach fixed point theorem}

% Sonnet 5 (Medium)
\section{Contraction mappings and the Banach fixed point theorem}
\label{sec:banach_fixed_point}

% Not in Corsin's own notes either, but needed to justify the appeal to the "Banach Fixed Point
% Theorem" in the solution to ex:ai_contraction_cos below, and a standard companion to
% completeness (def:complete_metric_space), which is why it sits in this chapter.
% Transition: Claude Opus 5 (High)
Completeness guarantees that certain limits exist, but on its own it does not say what they are.
The following theorem is the standard way of turning that guarantee into a construction: it
identifies a single point, proves it exists, proves it is unique, and supplies an algorithm that
converges to it. Almost every existence proof in the second half of this course goes through it.

\begin{definition}[Contraction mapping]
\label{def:contraction_mapping}
Let $(X,d)$ be a metric space. A map $f : X \to X$ is a \newterm{contraction}
\germanfor{Kontraktion} if it is Lipschitz continuous (\cref{def:lipschitz_continuous}) with
some Lipschitz constant $L \in [0,1)$, i.e.,
\[d(f(x),f(y)) \leq L \cdot d(x,y) \quad \text{for all } x,y \in X.\]
\end{definition}

\begin{theorem}[Banach fixed point theorem]
\label{thm:banach_fixed_point}
Let $(X,d)$ be a nonempty complete metric space (\cref{def:complete_metric_space}) and let
$f : X \to X$ be a contraction (\cref{def:contraction_mapping}). Then $f$ has a unique
\newterm{fixed point} \germanfor{Fixpunkt} $x^* \in X$, i.e., $f(x^*) = x^*$; moreover, for any
$x_0 \in X$, the iterates $x_{n+1} := f(x_n)$ converge to $x^*$.
\end{theorem}

% Supplement: Sascha Brack/Class Notes/Week_02_Notes_Monday_Updated.pdf, p. 13
\begin{proof}[Proof idea]
The core idea is to create a Cauchy sequence by repeatedly applying the contraction mapping $f$.
Fix an arbitrary starting point $x_0 \in X$, and define the sequence by $x_{n+1} := f(x_n)$ for all $n \geq 0$.
Comparing consecutive terms gives:
\[
d(x_{n+1}, x_n) = d(f(x_n), f(x_{n-1})) \leq \lambda \cdot d(x_n, x_{n-1}).
\]
Iterating this inequality yields $d(x_{n+1}, x_n) \leq \lambda^n d(x_1, x_0)$. Since $\lambda < 1$, one can use the triangle inequality and the geometric series to show that the distance between any two terms $x_m$ and $x_n$ (for $m > n$) becomes arbitrarily small as $n \to \infty$. Thus, the sequence is Cauchy.
Because $X$ is a complete metric space, this sequence converges to some limit $x^* \in X$. Since contraction mappings are continuous, one can pass the limit inside the function to conclude $f(x^*) = x^*$.
\end{proof}

% Generator: Claude Opus 5 (High) --- Sascha's proof idea establishes existence only; uniqueness
% is two lines and is half of what the theorem claims.
\begin{proof}[Uniqueness of the fixed point]
Suppose $f(x) = x$ and $f(y) = y$. Then
\[d(x,y) = d\big(f(x),f(y)\big) \leq L\,d(x,y),\]
so $(1-L)\,d(x,y) \leq 0$. Since $L<1$ we have $1-L>0$, forcing $d(x,y) \leq 0$ and hence
$d(x,y)=0$, i.e.\ $x=y$.
\end{proof}

% Generator: Claude Opus 5 (High) --- FIG-06-01, cobweb diagrams for the fixed point iteration.
% This chapter had no figure at all, and the theorem's proof is a construction (iterate f from
% an arbitrary starting point) whose whole content is visual in one variable.
% Arithmetic check. LEFT: f(x) = 0.45x + 0.35, a contraction with L = 0.45 < 1. Fixed point
% x* = 0.35/(1-0.45) = 0.35/0.55 = 0.6364. From x0 = 0.05 the iterates are
%   0.05 -> 0.3725 -> 0.51763 -> 0.58293 -> 0.61232 -> 0.62554, climbing to x*.
% RIGHT: f(x) = x + 0.12, slope exactly 1, so L = 1 and NOT a contraction. The graph is
% parallel to the diagonal and never meets it, so there is no fixed point at all. From
% x0 = 0.10 the iterates are 0.22, 0.34, 0.46, 0.58, 0.70, 0.82, marching off to the right.
% This is the second bullet of rem:banach_hypotheses (the T(x) = x+1 example) drawn on [0,1].
\begin{center}
\begin{tikzpicture}[x={(3.4cm,0)}, y={(0,3.4cm)}, >=Latex,
    web/.style={BrickRed, thick},
    fgraph/.style={blue!65!black, very thick}]

  % ---------------- left: a genuine contraction ----------------
  \begin{scope}
    \draw[->, gray!70] (0,0) -- (1.1,0) node[right, font=\footnotesize, text=black] {$x$};
    \draw[->, gray!70] (0,0) -- (0,1.1);
    \draw[gray!55, dashed] (0,0) -- (1,1) node[above left, font=\scriptsize, text=black] {$y=x$};
    \draw[fgraph] (0,0.35) -- (1,0.80) node[right, font=\scriptsize] {$f$};

    \draw[web] (0.05,0.05) -- (0.05,0.3725) -- (0.3725,0.3725) -- (0.3725,0.51763)
      -- (0.51763,0.51763) -- (0.51763,0.58293) -- (0.58293,0.58293) -- (0.58293,0.61232)
      -- (0.61232,0.61232) -- (0.61232,0.62554);

    \fill[black] (0.6364,0.6364) circle (1.7pt);
    \node[font=\scriptsize, anchor=west] at (0.68,0.60) {$x^{*}$};
    \draw[gray!70] (0.05,0.02) -- (0.05,-0.02) node[below, font=\scriptsize, text=black] {$x_0$};
    \node[font=\scriptsize, align=center, anchor=north] at (0.5,-0.15)
      {contraction, $L=0.45$\\ iterates climb to the unique $x^{*}$};
  \end{scope}

  % ---------------- right: slope exactly 1, no fixed point ----------------
  \begin{scope}[shift={(1.42,0)}]
    \draw[->, gray!70] (0,0) -- (1.1,0) node[right, font=\footnotesize, text=black] {$x$};
    \draw[->, gray!70] (0,0) -- (0,1.1);
    \draw[gray!55, dashed] (0,0) -- (1,1) node[below right, font=\scriptsize, text=black] {$y=x$};
    \draw[fgraph] (0,0.12) -- (0.88,1.0) node[above left, font=\scriptsize] {$f$};

    \draw[web] (0.10,0.10) -- (0.10,0.22) -- (0.22,0.22) -- (0.22,0.34) -- (0.34,0.34)
      -- (0.34,0.46) -- (0.46,0.46) -- (0.46,0.58) -- (0.58,0.58) -- (0.58,0.70)
      -- (0.70,0.70) -- (0.70,0.82) -- (0.82,0.82) -- (0.82,0.94);
    \draw[web, ->] (0.82,0.94) -- (0.94,0.94);

    \draw[gray!70] (0.10,0.02) -- (0.10,-0.02) node[below, font=\scriptsize, text=black] {$x_0$};
    \node[font=\scriptsize, align=center, anchor=north] at (0.5,-0.15)
      {$L=1$: graph parallel to $y=x$\\ no fixed point, iterates escape};
  \end{scope}
\end{tikzpicture}
\end{center}

The left panel is the proof in miniature. Each step moves vertically to the graph of $f$ and then
horizontally back to the diagonal, and because the graph is shallower than the diagonal, the steps
shrink by a factor of at least $L$ every time. That is precisely the geometric-series estimate the
proof runs on. The right panel shows what a single missing hypothesis costs: with slope exactly
$1$ the graph never meets the diagonal, so no fixed point exists and the same iteration marches
off instead of converging.

\begin{remark}[Where each hypothesis is used]
\label{rem:banach_hypotheses}
Both hypotheses of \cref{thm:banach_fixed_point} are indispensable, and each fails in a
different way.
\begin{itemize}
  \item \textbf{Completeness} is what supplies the limit. On the incomplete space
    $X := (0,1]$, the map $f(x) := x/2$ is a contraction with $L=\tfrac12$ and has no fixed
    point in $X$, since the iterates march towards $0$, which is exactly the point missing.
  \item \textbf{$L<1$ strictly} is what forces the iterates together. The map
    $f(x) := x + \tfrac{1}{1+e^{x}}$ on the complete space $\mathbb{R}$ satisfies the strict
    inequality $|f(x)-f(y)| < |x-y|$ for $x \neq y$, yet has no fixed point, since
    $f(x)>x$ everywhere. A supremum of $1$ that is never attained is still not good enough.
\end{itemize}
The uniqueness proof above shows the second point cleanly: the whole argument is the division by
$1-L$, which is exactly what $L=1$ forbids.
\end{remark}

% Supplement: Analysis_II_Script_v1.pdf, p. 20
\begin{exercise}[Both hypotheses of Banach's theorem are sharp]
\label{ex:banach_hypotheses_sharp}
Find examples for:
\begin{enumerate}[label=\textbf{(\alph*)}]
  \item A Lipschitz contraction $T : X \to X$ on a non-complete metric space $X$ that has no
    fixed point.
  \item A complete metric space $(X,d)$ and an \newterm{isometry} (a mapping $T : X \to X$ with
    $d(T(x_1),T(x_2)) = d(x_1,x_2)$ for all $x_1,x_2$) that has no fixed point.
\end{enumerate}
\end{exercise}

\begin{remark}
Part \textbf{(b)} sharpens \cref{rem:banach_hypotheses}'s second bullet: that example already
showed a Lipschitz constant of exactly $1$ can defeat the theorem, but its map $T$ was not an
isometry ($|T(x)-T(y)| < |x-y|$ strictly, just with no uniform gap). An isometry is the most
rigid failure imaginable, not shrinking distances even a little, and it \emph{still} need not have
a fixed point, even on a complete space. Note that unboundedness is doing real work in the
counterexample: on a \emph{compact} space an isometry into itself is automatically surjective, and
even that is not enough to force a fixed point. An irrational rotation of the circle $S^1$, which
is compact and an isometry, has none either. Fixed points of isometries are a
genuinely separate question from \cref{thm:banach_fixed_point}, not one it settles as a special
case.
\end{remark}

\begin{aiexercise}[{Contraction Mapping on $[0,1]$}]
\label{ex:ai_contraction_cos}
% Generator: Gemini 3.6 Flash (Medium)
Let $f \colon [0,1] \to [0,1]$ be defined by $f(x) := \frac{1}{2}\cos(x)$.
\begin{enumerate}[label=\textbf{(\alph*)}]
  \item Show that $f$ is a contraction mapping (\cref{def:contraction_mapping}) with Lipschitz
    constant $L = \frac{1}{2}\sin(1) < 1$.
  \item Deduce, using \cref{thm:banach_fixed_point}, that $f$ has a unique fixed point
    $x^* \in [0,1]$.
\end{enumerate}
\end{aiexercise}

% Supplement: Linus Lüchinger/Slides/Slides-03-02.pdf, slide 16;
%             Linus Lüchinger/Slides/Slides-03-05.pdf, slide 3
\begin{exercise}[Does the fixed-point theorem apply?]
\label{ex:bfpt_applicability_table}
\cref{rem:banach_hypotheses} shows that each hypothesis is needed by exhibiting one failure
apiece. Linus drills the same point the other way round, as a checklist: for each row below,
decide whether \cref{thm:banach_fixed_point} \emph{applies} to $T : X \to X$, and separately
determine the set of fixed points of $T$ in $X$. The two questions have different answers more
often than one expects.
\begin{center}
\begin{tabular}{lll}
\textbf{$T(x)$} & \textbf{$X$} & \textbf{Applies? / fixed points} \\
\hline
$\sin(x)$ & $\mathbb{R}$ & \\
$\tfrac{9}{10}\cos(x)$ & $\mathbb{R}$ & \\
$x^2$ & $\mathbb{R}$ & \\
$\tfrac12x^2$ & $(-1,1)$ & \\
$\tfrac12\big(x + \tfrac2x\big)$ & $\mathbb{Q}\cap(1,2)$ &
\end{tabular}
\end{center}
\end{exercise}


% --- exercise relocated here from the dissolved sheet 3 ---

% Originally: content/exercise-sheets/week-03.tex, problem 3.5
% Source: exercises/Ex3_Analysis2_eng.pdf

\begin{exercise}[Multiple choice (fixed points)]
\label{ex:3.5}
Select all the statements below that are true.
\begin{enumerate}[label=\textbf{(\alph*)}]
  \item If you spread a map of Zurich on your desk, then one point of the desk will coincide with
    its representation on the map (in an ideal world).
  \item If $f : [0,1] \to [0,2]$ is $\tfrac{1}{2}$-Lipschitz, then $f$ has a fixed point.
  \item If $f : [0,2] \to [0,1]$ is $\tfrac{1}{2}$-Lipschitz, then $f$ has a fixed point.
  \item If $f : [0,1] \cup [2,3] \to [0,1] \cup [2,3]$ is continuously differentiable with
    $|f'(x)| < 1$ for all $x \in [0,1]\cup[2,3]$, then it has a fixed point.
  \item \textbf{(*)} If $f : [0,1] \to [0,1]$ is differentiable with $|f'(x)| < 1$ for all
    $x \in (0,1)$, then it has a fixed point.
\end{enumerate}
\exinfo{This exercise is Problem 3.5 of Problem Sheet 3, Analysis II, Spring Semester 2026.}
\end{exercise}

% Originally: the \section{Solutions} of content/week-02.tex and content/week-03.tex

\section{Solutions}
\label{sec:completeness_solutions}

\begin{exercisesolution}[Proof of \cref{ex:cauchy_elementary_facts}]
% Generator: Claude Sonnet 5 (High)
\textbf{(a) Bounded.} Take $\varepsilon=1$: there is $N$ with $d(x_n,x_m)<1$ for all
$n,m\geq N$; in particular $d(x_n,x_N)<1$ for $n\geq N$. Setting
$M := \max\{d(x_0,x_N),\dots,d(x_{N-1},x_N),\,1\}$, the triangle inequality gives
$d(x_n,x_0) \leq d(x_n,x_N)+d(x_N,x_0) \leq 1+M$ for every $n$, so $(x_n)_n \subseteq B_{1+M}(x_0)$.

\textbf{(b) Convergent $\Rightarrow$ Cauchy.} If $x_n \to x$, fix $\varepsilon>0$ and choose $N$
with $d(x_n,x)<\varepsilon/2$ for $n\geq N$. Then for $n,m\geq N$,
$d(x_n,x_m) \leq d(x_n,x)+d(x,x_m) < \varepsilon$.

\textbf{(c) Cauchy + convergent subsequence $\Rightarrow$ convergent.} Let $(x_n)$ be Cauchy and
$x_{n_k} \to x$. The \qt{only if} direction is part \textbf{(b)} applied to the whole sequence.
For \qt{if}: fix $\varepsilon>0$. Cauchyness gives $N$ with $d(x_n,x_m)<\varepsilon/2$ for
$n,m\geq N$; convergence of the subsequence gives $K$ with $d(x_{n_k},x)<\varepsilon/2$ for
$k\geq K$. Since $n_k \to \infty$, pick $k\geq K$ with $n_k \geq N$. Then for every $n\geq N$,
\[d(x_n,x) \leq d(x_n,x_{n_k}) + d(x_{n_k},x) < \tfrac{\varepsilon}{2}+\tfrac{\varepsilon}{2} = \varepsilon,\]
so $x_n \to x$.
\end{exercisesolution}

\begin{exercisesolution}[Solution to \cref{ex:completeness_not_topological}]
% Generator: Claude Opus 5 (Medium)
Write $\varphi : (0,1] \to [1,\infty)$, $\varphi(x) := 1/x$. It is a bijection, and by
construction $d(x,y) = |\varphi(x)-\varphi(y)|$. Consequently, $d$ is nothing but the Euclidean
metric of $[1,\infty)$ transported back to $(0,1]$ along $\varphi$. Every part follows from that.

\begin{enumerate}[label=\textbf{(\alph*)}]
  \item Symmetry and the triangle inequality are inherited from $|\cdot|$ directly. For
    definiteness, $d(x,y)=0$ forces $\varphi(x)=\varphi(y)$, and $\varphi$ is injective, so
    $x=y$.
  \item The sequence $x_n := \tfrac1n$ lies in $(0,1]$ and is Cauchy for
    $d_{\text{eucl}}$, being Cauchy in $\mathbb{R}$. Its limit in $\mathbb{R}$ is $0 \notin X$,
    and limits in $\mathbb{R}$ are unique (\cref{prop:uniqueness_of_limits}), so it has no limit
    in $X$.
  \item Let $(x_n)$ be $d$-Cauchy. Then $\big(\varphi(x_n)\big)$ is Cauchy in $\mathbb{R}$,
    hence converges to some $\ell \in \mathbb{R}$. Since $\varphi(x_n) \geq 1$ for every $n$,
    also $\ell \geq 1$; in particular $\ell \neq 0$, so $x := 1/\ell$ lies in $(0,1]$. Then
    \[d(x_n,x) = \big|\varphi(x_n) - \varphi(x)\big| = \big|\varphi(x_n)-\ell\big|
      \longrightarrow 0,\]
    so $x_n \to x$ in $(X,d)$. Note where the boundedness of $X$ entered: it is what forces
    $\ell \geq 1$ rather than merely $\ell \geq 0$.
  \item $\varphi$ and $\varphi^{-1}(t) = 1/t$ are continuous, so $\varphi$ is a
    homeomorphism. A set is $d$-open exactly when its $\varphi$-image is open in $[1,\infty)$,
    exactly when it is $d_{\text{eucl}}$-open. Concretely, the same two-sided-ball argument as in
    \Cref{ex:heine_borel_fails}\textbf{(b)} works: on $(0,1]$ one has
    $|x-y| \leq d(x,y) \leq \tfrac{1}{x y}\,|x-y|$ with $\tfrac{1}{xy}$ locally bounded.
\end{enumerate}

\begin{remark}
So $d_{\text{eucl}}$ and $d$ have the \emph{same open sets}, the same convergent sequences, the
same continuous functions and the same compact sets, and yet one is complete and the other is
not. \textbf{Completeness is a property of the metric, not of the topology it generates.}

This completes a trio worth holding together, all three proved by the same device of comparing
two metrics that induce one topology:
\begin{itemize}
  \item \textbf{Compactness is topological.} It is defined by open covers alone
    (\cref{def:compact_metric_space}), so it cannot distinguish two metrics with the same open
    sets.
  \item \textbf{Boundedness is not:} see \cref{rem:boundedness_not_topological}.
  \item \textbf{Completeness is not:} this exercise.
\end{itemize}
That is exactly why \qt{compact if and only if closed and bounded} needs the standard metric, and
why the correct general replacement is \qt{complete and totally bounded}
(\cref{rem:compact_iff_complete_totally_bounded}). Both ingredients of the replacement are
non-topological, and only their combination is.
\end{remark}
\end{exercisesolution}

\begin{exercisesolution}[Proof of \cref{ex:completion_well_defined}]
% Generator: Claude Sonnet 5 (High)
\textbf{$\sim$ is an equivalence relation on $\mathcal{C}_X$.} Reflexivity is
$d(x_n,x_n) = 0 \to 0$. Symmetry is symmetry of $d$. For transitivity, if
$(x_n)\sim(y_n)$ and $(y_n)\sim(z_n)$, the triangle inequality gives
$0 \leq d(x_n,z_n) \leq d(x_n,y_n)+d(y_n,z_n) \to 0+0 = 0$, so $(x_n)\sim(z_n)$.

\textbf{$\overline d$ is well-defined.} First, the limit defining $\overline d$ exists: for
$(x_n),(y_n) \in \mathcal{C}_X$, the reverse triangle inequality gives
\[\big|d(x_n,y_n) - d(x_m,y_m)\big| \leq d(x_n,x_m) + d(y_n,y_m),\]
and the right-hand side $\to 0$ as $n,m\to\infty$ since $(x_n)$ and $(y_n)$ are each Cauchy, so
$\big(d(x_n,y_n)\big)_n$ is a Cauchy sequence of \emph{real numbers}, hence convergent. Note that
this is the completeness axiom for $\mathbb{R}$ itself from \emph{Analysis I}, and not an instance
of \cref{def:complete_metric_space} being proved circularly. Second, the limit
does not depend on the choice of representative: if $(x_n)\sim(x_n')$ and $(y_n)\sim(y_n')$, the
same reverse triangle inequality gives
$\big|d(x_n,y_n)-d(x_n',y_n')\big| \leq d(x_n,x_n')+d(y_n,y_n') \to 0$, so the two limits agree.

\textbf{$\overline d$ is a metric.} Fix representatives $(x_n)\in\overline x$, $(y_n)\in\overline
y$, $(z_n)\in\overline z$.
\begin{itemize}
  \item \textbf{Definiteness.} $\overline d(\overline x,\overline y) = 0$ means
    $\lim_n d(x_n,y_n) = 0$, i.e.\ exactly $(x_n)\sim(y_n)$, i.e.\ $\overline x = \overline y$.
  \item \textbf{Symmetry.} Immediate from $d(x_n,y_n)=d(y_n,x_n)$ for every $n$.
  \item \textbf{Triangle inequality.} Take $n\to\infty$ in
    $d(x_n,z_n) \leq d(x_n,y_n) + d(y_n,z_n)$, using that limits preserve weak inequalities.
\end{itemize}
Finally $d(x,y) = \overline d(\iota(x),\iota(y))$ holds because the constant sequences
$(x,x,\dots)$ and $(y,y,\dots)$ give $d(x_n,y_n) = d(x,y)$ for every $n$, so the limit is
$d(x,y)$ itself; and this identity makes $\iota$ injective, since $d(x,y)=0 \iff x=y$.
\end{exercisesolution}

\begin{exercisesolution}[Proof of \cref{ex:completion_is_complete}]
% Generator: Claude Sonnet 5 (High)
Let $(\xi_m)_{m=0}^\infty$ be Cauchy in $\overline X$, with $\xi_m = [(x_{m,n})_{n=0}^\infty]$ for
some Cauchy sequence $(x_{m,n})_n$ in $X$ representing $\xi_m$. Replacing $(x_{m,n})_n$ by a tail
of itself if necessary, which is harmless because a Cauchy sequence and any of its tails represent
the same class, differing in only finitely many terms, we may assume each representative is already
\newterm{regularized}:
\[d(x_{m,n},x_{m,n'}) < \tfrac1m \qquad \text{for all } n,n' \geq 0.\]
Set $y_m := x_{m,m}$, the diagonal sequence the exercise asks about. Letting $n'\to\infty$ in the
regularization bound (with $n=m$) gives
\begin{equation}
\overline d\big(\iota(y_m),\,\xi_m\big) = \lim_{n\to\infty} d(x_{m,m},x_{m,n}) \ \leq\ \tfrac1m.
\label{eq:completion_diagonal_close}
\end{equation}

\textbf{$(y_m)_m$ is Cauchy in $X$.} By the triangle inequality for $\overline d$ (already
established in \cref{ex:completion_well_defined}) and \eqref{eq:completion_diagonal_close},
\[d(y_m,y_{m'}) = \overline d\big(\iota(y_m),\iota(y_{m'})\big)
  \ \leq\ \overline d(\iota(y_m),\xi_m) + \overline d(\xi_m,\xi_{m'}) + \overline d(\xi_{m'},\iota(y_{m'}))
  \ \leq\ \tfrac1m + \overline d(\xi_m,\xi_{m'}) + \tfrac1{m'}.\]
Given $\varepsilon>0$, Cauchyness of $(\xi_m)_m$ gives $M_1$ with
$\overline d(\xi_m,\xi_{m'}) < \varepsilon/3$ for $m,m'\geq M_1$; taking also
$M \geq \max\{M_1,\,3/\varepsilon\}$ makes all three terms $< \varepsilon/3$ for $m,m'\geq M$, so
$(y_m)_m$ is Cauchy in $X$ and defines a class $\overline x := [(y_m)_m] \in \overline X$.

\textbf{$\xi_m \to \overline x$.} Using \eqref{eq:completion_diagonal_close} again and the same
triangle inequality,
\[\overline d(\xi_m,\overline x) \ \leq\ \overline d(\xi_m,\iota(y_m)) + \overline d(\iota(y_m),\overline x)
  \ \leq\ \tfrac1m + \lim_{k\to\infty}d(y_m,y_k).\]
The first term $\to 0$ as $m\to\infty$ by construction, and the second does too, being (for each
fixed $m$) bounded by $\sup_{k\geq m}d(y_m,y_k)$, which $\to0$ as $m\to\infty$ since $(y_m)_m$ is
Cauchy. Hence $\overline d(\xi_m,\overline x) \to 0$, i.e.\ $\xi_m \to \overline x$ in
$\overline X$, as required.
\end{exercisesolution}

\begin{exercisesolution}[Proof of \cref{ex:completion_universal_property}]
% Generator: Claude Sonnet 5 (High)
\textbf{Construction.} For $\overline x = [(x_n)_n] \in \overline X$, the sequence $(f(x_n))_n$
is Cauchy in $Y$, since $d_Y(f(x_n),f(x_m)) = d(x_n,x_m) \to 0$ as $(x_n)_n$ is Cauchy in $X$.
As $Y$ is complete, $f(x_n) \to y$ for a unique $y \in Y$
(\cref{prop:uniqueness_of_limits}); define $\overline f(\overline x) := y$.

\textbf{Well-defined.} If $(x_n)\sim(x_n')$, then $d_Y(f(x_n),f(x_n')) = d(x_n,x_n') \to 0$, so
$(f(x_n))_n$ and $(f(x_n'))_n$ have the same limit in $Y$ by uniqueness of limits; hence
$\overline f(\overline x)$ does not depend on the representative chosen.

\textbf{$f = \overline f \circ \iota_X$.} For $x\in X$, $\iota_X(x)$ is represented by the
constant sequence $(x,x,\dots)$, so $\overline f(\iota_X(x)) = \lim_n f(x) = f(x)$.

\textbf{$\overline f$ satisfies the displayed identity.} For $\overline x = [(x_n)_n]$,
$\overline y = [(y_n)_n]$, continuity of $d_Y$ in both arguments (itself $1$-Lipschitz, by the
same reverse-triangle-inequality argument as in \cref{ex:completion_well_defined}) gives
\[d_Y\big(\overline f(\overline x),\overline f(\overline y)\big)
  = d_Y\Big(\lim_n f(x_n),\ \lim_n f(y_n)\Big)
  = \lim_n d_Y\big(f(x_n),f(y_n)\big)
  = \lim_n d(x_n,y_n) = \overline d(\overline x,\overline y).\]

\textbf{Uniqueness.} Let $g : \overline X \to Y$ satisfy the same two conditions. The displayed
identity makes $g$ (like $\overline f$) $1$-Lipschitz, hence continuous. For
$\overline x = [(x_n)_n]$, the sequence $\iota_X(x_n)$ converges to $\overline x$ in
$\overline X$: indeed $\overline d(\iota_X(x_n),\overline x) = \lim_k d(x_n,x_k) \to 0$ as
$n\to\infty$, since $(x_n)_n$ is Cauchy. By continuity of $g$ and $f = g\circ\iota_X$,
\[g(\overline x) = \lim_n g(\iota_X(x_n)) = \lim_n f(x_n) = \overline f(\overline x),\]
using uniqueness of limits in $Y$ once more for the last equality. As $\overline x$ was
arbitrary, $g = \overline f$.
\end{exercisesolution}

\begin{exercisesolution}[Proof of \cref{ex:banach_hypotheses_sharp}]
% Generator: Claude Sonnet 5 (High)
\begin{enumerate}[label=\textbf{(\alph*)}]
  \item This is \cref{rem:banach_hypotheses}'s first bullet: $X := (0,1]$ is not complete,
    and $T(x) := x/2$ is a $\tfrac12$-contraction on $X$ with iterates marching towards the
    missing point $0$, hence no fixed point in $X$.
  \item Take $X := \mathbb{R}$ (complete, \cref{prop:Rn_coordinatewise_complete}) with the
    translation $T(x) := x+1$. Then $d(T(x),T(y)) = |(x+1)-(y+1)| = |x-y| = d(x,y)$, so $T$ is an
    isometry, and in particular Lipschitz with constant $L=1$ rather than $L<1$. Consequently,
    this does not contradict \cref{thm:banach_fixed_point}. A fixed point would need $x+1=x$,
    which is impossible, so $T$ has none.
\end{enumerate}
\end{exercisesolution}

\begin{exercisesolution}[Proof of \cref{ex:kobel_local_banach_fixed_point}]
% Generator: Claude Opus 5 (High)
Write $K := \overline{B}_r(x_0)$. The whole content of the exercise is that $f$ restricts to a
self-map of $K$, after which \cref{thm:banach_fixed_point} finishes the job.

\textbf{$K$ is complete.} A closed ball is a closed subset of $X$, and a closed subset of a
complete metric space is complete: a Cauchy sequence in $K$ is Cauchy in $X$, hence converges in
$X$, and its limit lies in $K$ because $K$ is closed. It is also non-empty, since $x_0 \in K$.

\textbf{$f(K) \subseteq K$.} Let $x \in K$, so $d(x,x_0) \le r$. Since $K \subseteq E$, both $x$
and $x_0$ lie in the domain of $f$, and the triangle inequality together with the contraction
estimate gives
\[d\big(f(x),x_0\big) \ \le\ d\big(f(x),f(x_0)\big) + d\big(f(x_0),x_0\big)
  \ \le\ L\,d(x,x_0) + (1-L)r \ \le\ Lr + (1-L)r \ =\ r,\]
so $f(x) \in K$. The hypothesis $d(f(x_0),x_0) \le (1-L)r$ is used exactly once, here, and it is
what makes the two terms add up to $r$ rather than to something larger.

\textbf{Conclusion.} The restriction $f|_K : K \to K$ is a contraction with the same constant $L$
on the non-empty complete metric space $K$, so \cref{thm:banach_fixed_point} supplies a unique
$\bar x \in K$ with $f(\bar x) = \bar x$. Since $K \subseteq E$, this $\bar x$ is a fixed point in
$E$, as required.

\textbf{Uniqueness comes for free, and in all of $E$.} \Cref{thm:banach_fixed_point} gives
uniqueness only within $K$, but the two-line argument used for the theorem itself needs nothing
beyond the contraction estimate, which holds on all of $E$: if $f(x)=x$ and $f(y)=y$ with
$x,y \in E$, then $d(x,y) = d(f(x),f(y)) \le L\,d(x,y)$, and $L<1$ forces $d(x,y)=0$. So $\bar x$
is the only fixed point of $f$ in $E$, even though the existence argument saw only the ball.
\end{exercisesolution}

\begin{exercisesolution}[Proof of \cref{ex:ai_contraction_cos}]
% Generator: Gemini 3.6 Flash (Medium)
\begin{enumerate}[label=\textbf{(\alph*)}]
  \item By the Mean Value Theorem, for any $x, y \in [0,1]$, there exists $\xi \in (0,1)$ such that $|f(x) - f(y)| = |f'(\xi)| |x-y| = \frac{1}{2}\sin(\xi) |x-y|$. Since $\sin(x)$ is strictly increasing on $[0,1]$, $\sup_{\xi \in [0,1]} \frac{1}{2}\sin(\xi) = \frac{1}{2}\sin(1) \approx 0.4207 < 1$. Thus $f$ is a contraction (\cref{def:contraction_mapping}) with $L = \frac{1}{2}\sin(1) < 1$.
  \item Since $[0,1]$ equipped with the Euclidean metric is a complete metric space (\cref{def:complete_metric_space}) and $f([0,1]) \subseteq [0,1]$, \cref{thm:banach_fixed_point} guarantees the existence and uniqueness of a fixed point $x^* \in [0,1]$ satisfying $f(x^*) = x^*$.
\end{enumerate}
\end{exercisesolution}

\begin{exercisesolution}[Solution to \cref{ex:3.5}]
% Generator: Claude Opus 5 (High)
\textbf{(a), (c) and (e) are true; (b) and (d) are false.}

\begin{enumerate}[label=\textbf{(\alph*)}]
  \item \textbf{True,} and this is the classic illustration of \cref{thm:banach_fixed_point}. The
    map of Zurich lying on the desk defines a function from the (compact, hence complete) region
    of the city depicted onto a smaller region of the desk, and that function is a similarity: a
    rotation and translation composed with a scaling by the map's ratio $L < 1$. It is therefore
    a contraction of a complete metric space into itself, and so has exactly one fixed point: one
    spot on the map lying directly above the point of the city it represents.
  \item \textbf{False.} The hypothesis says nothing about $f$ mapping its domain into itself, and
    the codomain $[0,2]$ is strictly larger than the domain $[0,1]$. Take
    \[f(x) := \tfrac{x}{2} + \tfrac32, \qquad f : [0,1] \to [\tfrac32,2] \subseteq [0,2],\]
    which is $\tfrac12$-Lipschitz. A fixed point would need $x = \tfrac x2 + \tfrac32$, i.e.\
    $x=3$, which is not in $[0,1]$.
  \item \textbf{True.} Now the inclusion runs the right way: $f$ maps $[0,2]$ into
    $[0,1] \subseteq [0,2]$, so $f$ is a self-map of the complete space $[0,2]$
    (closed and bounded in $\mathbb{R}$, hence compact, hence complete by
    \cref{ex:compact_finite_complete}), and it is a contraction with $L = \tfrac12 < 1$.
    \Cref{thm:banach_fixed_point} applies directly.
  \item \textbf{False,} and this is the instructive one: the domain is \emph{disconnected}, and
    the condition $|f'| < 1$ is purely local, so it cannot see across the gap. Define
    \[f(x) := \begin{cases} \tfrac52 & x \in [0,1], \\ \tfrac12 & x \in [2,3].\end{cases}\]
    Then $f$ is continuously differentiable with $f' \equiv 0$, it maps
    $[0,1]\cup[2,3]$ into itself, and it has no fixed point, because each of the two pieces is
    sent into the \emph{other} one. Compare the local-to-global reading of
    \cref{rem:local_to_global}: a pointwise derivative bound upgrades to a global contraction
    only along paths that stay inside the domain, and here there are none.
  \item \textbf{True,} but \emph{not} by the fixed point theorem, and that is the point of the
    star. The hypothesis $|f'(x)| < 1$ gives no uniform $L < 1$, since the supremum of $|f'|$ may
    well equal $1$, so \cref{thm:banach_fixed_point} is unavailable. Existence follows from the
    intermediate value theorem instead. Put $g(x) := f(x) - x$, which is continuous on $[0,1]$.
    Since $f$ takes values in $[0,1]$,
    \[g(0) = f(0) \geq 0 \qquad\text{and}\qquad g(1) = f(1) - 1 \leq 0,\]
    so $g$ vanishes somewhere in $[0,1]$ by \cref{cor:intermediate_value_theorem}, and that zero
    is a fixed point of $f$. Note that the derivative hypothesis was never used: \emph{every}
    continuous self-map of $[0,1]$ has a fixed point. What $|f'| < 1$ buys is
    \textbf{uniqueness}, by the same two-line argument as in the proof of uniqueness for
    \cref{thm:banach_fixed_point} combined with the mean value theorem.
\end{enumerate}
\end{exercisesolution}

\begin{exercisesolution}[Solution to \cref{ex:bfpt_applicability_table}]
% Generator: Claude Opus 5 (High)
Recall that \cref{thm:banach_fixed_point} needs three things at once: $X$ complete, $T$ mapping
$X$ into $X$, and a Lipschitz constant $\lambda < 1$ that is \emph{actually attained as a bound},
not merely approached.
\begin{itemize}
  \item \textbf{$T(x) = \sin(x)$ on $\mathbb{R}$. Does not apply; fixed points $\{0\}$.}
    $\mathbb{R}$ is complete and $T$ maps it into itself, but
    $\lambda = \sup_{x}\lvert\cos x\rvert = 1$, attained at $x=0$. So the theorem is unavailable.
    A fixed point exists all the same: $\sin x = x$ holds at $x=0$ and nowhere else, since
    $\lvert\sin x\rvert < \lvert x\rvert$ for $x \neq 0$. This is the first lesson of the table:
    the theorem is \emph{sufficient}, never necessary.
  \item \textbf{$T(x) = \tfrac{9}{10}\cos(x)$ on $\mathbb{R}$. Applies; one fixed point.}
    Here $\lvert T'(x)\rvert = \tfrac{9}{10}\lvert\sin x\rvert \leq \tfrac{9}{10} < 1$, so by the
    mean value theorem $T$ is a $\tfrac{9}{10}$-contraction; $\mathbb{R}$ is complete and
    $T(\mathbb{R}) \subseteq [-\tfrac9{10},\tfrac9{10}] \subseteq \mathbb{R}$. All three
    hypotheses hold, so there is exactly one fixed point. It has no closed form; iterating from
    any starting value gives $x^* \approx 0.693$.
  \item \textbf{$T(x) = x^2$ on $\mathbb{R}$. Does not apply; fixed points $\{0,1\}$.}
    $T$ is not Lipschitz at all on $\mathbb{R}$, since $T'(x) = 2x$ is unbounded. That it cannot
    apply is visible without any computation: $x^2 = x$ has the two solutions $0$ and $1$, and
    the theorem asserts \emph{uniqueness}, so any set of hypotheses implying it must fail here.
  \item \textbf{$T(x) = \tfrac12x^2$ on $(-1,1)$. Does not apply; fixed points $\{0\}$.}
    Two hypotheses fail together. The domain $(-1,1)$ is not complete, and
    $\lambda = \sup_{\lvert x\rvert<1}\lvert T'(x)\rvert = \sup\lvert x\rvert = 1$. This supremum
    is not attained, but a supremum of $1$ is still not $<1$, which is exactly the trap of
    \cref{rem:banach_hypotheses}. The map does send $(-1,1)$ into $[0,\tfrac12) \subseteq
    (-1,1)$. Solving $\tfrac12x^2 = x$ gives $x \in \{0,2\}$, of which only $0$ lies in the
    domain, so a unique fixed point exists regardless.
  \item \textbf{$T(x) = \tfrac12\big(x+\tfrac2x\big)$ on $\mathbb{Q}\cap(1,2)$. Does not apply;
    \emph{no} fixed points.} This row is the interesting one, and it is the Babylonian iteration
    for $\sqrt2$. Everything works except completeness:
    \begin{itemize}
      \item $T'(x) = \tfrac12 - \tfrac1{x^2}$, so on $(1,2)$ we have
        $-\tfrac12 \leq T'(x) \leq \tfrac14$ and hence $\lambda = \tfrac12 < 1$: a genuine
        contraction.
      \item $T$ maps $X$ into $X$: it sends rationals to rationals, and on $(1,2)$ its minimum
        is $T(\sqrt2) = \sqrt2 \approx 1.414$ while $T(1) = T(2) = \tfrac32$, so
        $T\big((1,2)\big) \subseteq [\sqrt2,\tfrac32] \subseteq (1,2)$.
      \item But $\mathbb{Q}\cap(1,2)$ is \textbf{not complete}.
    \end{itemize}
    And there is indeed no fixed point: $x = \tfrac12(x+\tfrac2x)$ forces $x^2 = 2$, so
    $x = \sqrt2 \notin \mathbb{Q}$. The iterates are the familiar rational approximations
    $\tfrac32, \tfrac{17}{12}, \tfrac{577}{408}, \dots$, which form a Cauchy sequence in $X$
    converging to a point that $X$ does not contain.

    This is worth contrasting with \cref{rem:banach_hypotheses}, where completeness failed on
    $(0,1]$ with the missing point sitting at an endpoint one is already suspicious of. Here the
    domain is an interval of rationals, the hole is invisible, and the contraction is a
    thoroughly practical algorithm. Completeness is not a technicality.
\end{itemize}
\end{exercisesolution}

\begin{exercisesolution}[Solution of \cref{ex:incomplete_metric_space_example}]
% Generator: Claude Opus 5 (High)
Take $X := (0,1]$ with the Euclidean distance $d(x,y) := \lvert x-y\rvert$ inherited from
$\mathbb{R}$. The sequence $x_k := 1/k$ lies in $X$ and is Cauchy, since $\lvert x_k - x_l\rvert \le
\max\{1/k, 1/l\} \to 0$. But it does not converge in $X$: its only candidate limit in $\mathbb{R}$
is $0$, and $0 \notin X$, so no point of $X$ can be its limit. Hence $(X,d)$ is not complete.

The same mechanism gives the other standard answers. Taking $X := \mathbb{Q}$ with
$d(x,y) := \lvert x-y\rvert$, any rational sequence converging to $\sqrt2$ is Cauchy without a limit
in $X$; taking $X := \mathbb{R} \setminus \{0\}$ works for the same reason as $(0,1]$.
\begin{ainote}
What is being exploited each time is that completeness is \emph{not} a property of the distance
alone: $(0,1]$ and $[0,1]$ carry the same formula for $d$, and one is complete while the other is
not. A Cauchy sequence is a sequence that is trying to converge, and the question is only whether
the space has kept the point it is aiming at. This is exactly the gap that the completion of a
metric space (\cref{sec:completion_of_a_metric_space}) is built to fill.
\end{ainote}
\end{exercisesolution}

\begin{exercisesolution}[Solution of \cref{ex:kobel_cauchy_single_choice}]
% Generator: Claude Opus 5 (High)
Only \textbf{(d)} is correct.
\begin{enumerate}[label=\textbf{(\alph*)}]
  \item \textbf{False.} A Cauchy sequence converges precisely when the ambient space has kept the
    point it is aiming at, and that is the definition of completeness
    (\cref{def:complete_metric_space}), not a theorem about Cauchy sequences.
    \Cref{ex:incomplete_metric_space_example} is a counterexample in one line.
  \item \textbf{False}, and this is the option the question is built around. Closedness is not a
    property a space has, only a property a subset has \emph{relative to} something larger: every
    metric space is closed in itself, so reading (b) as a hypothesis on $X$ alone makes it say
    nothing at all, and $X := (0,1]$ is again a counterexample. What is true is the relative
    statement: a closed subset of a \emph{complete} space is complete, and the ambient
    completeness is the part (b) omits.
  \item \textbf{False.} An interval need not be complete either. On $X := (0,1)$ the sequence
    $a_i := 1/i$ is Cauchy with no limit in $X$. Only the closed bounded intervals, and the closed
    unbounded ones, are complete.
  \item \textbf{True.} A compact metric space is sequentially compact
    (\cref{thm:sequential_compactness}), so $(a_i)$ has a convergent subsequence, and a Cauchy
    sequence with a convergent subsequence converges
    (\cref{ex:cauchy_elementary_facts}\textbf{(c)}). Hence every compact metric space is complete,
    which is \cref{ex:compact_finite_complete}.
\end{enumerate}
\end{exercisesolution}

\begin{exercisesolution}[Solution of \cref{ex:banach_hypotheses_true_false}]
% Generator: Claude Opus 5 (High)
\begin{enumerate}[label=\textbf{(\alph*)}]
  \item \textbf{True.} For $x,y \in [0,1]$,
    $\lvert x^2 - y^2\rvert = \lvert x+y\rvert\,\lvert x-y\rvert \le 2\lvert x-y\rvert$,
    since $x + y \le 2$ on $[0,1]$. So $2$ is a Lipschitz constant. It is even the optimal one, as
    $\lvert f'(x)\rvert = 2x$ approaches $2$ at $x = 1$.
  \item \textbf{False.} The space $X = [0,1]$ is complete, and $f$ does map $X$ into $X$; what fails
    is the contraction hypothesis. A contraction needs a Lipschitz constant \emph{strictly less
    than $1$}, and by \textbf{(a)} the best available constant here is $2$. Concretely, taking
    $x = 1$ and $y = 1 - \varepsilon$ gives
    $\lvert f(x) - f(y)\rvert = (2-\varepsilon)\varepsilon > \lvert x - y\rvert$ for small
    $\varepsilon > 0$: near $x = 1$ the map \emph{expands} distances.
  \item \textbf{True} --- and note the plural. Solving $x^2 = x$ on $[0,1]$ gives $x = 0$ and
    $x = 1$, so $f$ has exactly two fixed points.
\end{enumerate}
\begin{remark}
Parts \textbf{(b)} and \textbf{(c)} together are the point of the exercise. The Banach fixed point
theorem (\cref{thm:banach_fixed_point}) is a sufficient condition, not a necessary one: its
hypotheses fail here, and fixed points exist anyway. What the failure does cost is the
\emph{uniqueness} half of the conclusion, and sure enough there are two of them. Reading \qt{the
hypotheses fail} as \qt{there is no fixed point} is the standard mistake.
\end{remark}
\end{exercisesolution}

\begin{exercisesolution}[Proof of \cref{ex:banach_fixed_point_linear}]
% Generator: Gemini 3.6 Flash (High)
% Correction: Claude Opus 5 (High) --- rewritten for R^n. Flash's version worked in an abstract
% metric space and wrote d(f(x)/lambda, f(y)/lambda) = d(f(x),f(y))/lambda, which is meaningless
% there and false for a general metric even when scaling does make sense. It is the homogeneity of
% the Euclidean norm that is doing the work, so the norm has to be visible in the computation.
% Confirmed against WS22_Lösung.pdf p. 16, which runs the identical norm computation and likewise
% names the threshold lambda > L. The official solution agrees line for line.
\begin{enumerate}[label=\textbf{(\alph*)}]
  \item Fix $x,y \in \mathbb{R}^n$. Since the Euclidean norm is homogeneous and $f$ is
    $L$-Lipschitz,
\[
\lvert g(x) - g(y)\rvert = \left\lvert \frac{f(x) - f(y)}{\lambda} \right\rvert
  = \frac{1}{\lambda}\,\lvert f(x) - f(y)\rvert \leq \frac{L}{\lambda}\,\lvert x - y\rvert .
\]
Set $k := L/\lambda$. Since $\lambda > L > 0$ we have $k < 1$, and therefore $g$ is a contraction
with contraction constant $k$.

  \item The space $\mathbb{R}^n$ is complete, so the Banach Fixed Point Theorem
    (\cref{thm:banach_fixed_point}) applies to $g$ and yields a unique $x \in \mathbb{R}^n$ with
    $g(x) = x$. Multiplying $\lambda^{-1} f(x) = x$ by $\lambda$ turns this into $f(x) = \lambda x$,
    and the two equations have exactly the same solutions. Hence $f(x) = \lambda x$ has a unique
    solution.
\end{enumerate}
\begin{ainote}
The hypothesis is what makes this work, and it is worth naming: a Lipschitz $f$ need not be a
contraction, and need have no fixed point at all (translation $x \mapsto x + v$ is $1$-Lipschitz).
Dividing by $\lambda > L$ buys the missing factor. This is also why the original exam question says
\qt{for all sufficiently large $\lambda$} rather than naming a threshold: any $\lambda > L$ works,
but $L$ is not part of the data the question hands you.
\end{ainote}
\end{exercisesolution}

\begin{exercisesolution}[Proof of \cref{ex:examHS24_TF_banach}]
% Generator: Gemini 3.1 Pro (High)
\begin{enumerate}[label=\textbf{(\alph*)}]
  \item \textbf{True:} The derivative is $f'(x) = -\sin(x)$. Since $\lvert f'(x)\rvert = \lvert\sin x\rvert \le 1$ for all $x \in \mathbb{R}$, the mean value theorem gives $\lvert f(x) - f(y)\rvert \le 1 \cdot \lvert x-y\rvert$, so $f$ has Lipschitz constant $1$.
  \item \textbf{False:} The Banach fixed point theorem requires a \emph{strict} contraction, that is, a Lipschitz constant $\lambda < 1$. Here, the tightest possible constant is $\sup_{x \in [-\pi,\pi]} \lvert f'(x)\rvert = 1$ (attained at $x = \pm \pi/2$), so the map is non-expansive but not strictly contractive.
  \item \textbf{True:} The function $f(x) = \cos x$ is continuous and maps $[-\pi,\pi]$ into $[-1,1] \subseteq [-\pi,\pi]$. Since $g(x) := \cos x - x$ has $g(-\pi) = \pi - 1 > 0$ and $g(\pi) = -1 - \pi < 0$, the intermediate value theorem supplies a root, which is a fixed point of $f$.
\end{enumerate}
\end{exercisesolution}



